\chapter{并发重访}
\label{CH:LOCK2}

在内核设计中，同时实现良好的并行性能、确保并发正确性并保持代码可理解性，是一个巨大的挑战。直接使用锁是实现正确性的最佳途径，但并非总是可行。本章重点介绍 xv6 被迫以复杂方式使用锁的例子，以及 xv6 使用类锁技术但不用锁的例子。

\section{锁模式}

缓存项通常对加锁是一个挑战。例如，文件系统的块缓存 \lineref{kernel/bio.c:/^struct/} 存储最多 {\tt NBUF} 个磁盘块的副本。给定磁盘块在缓存中最多只有一个副本是至关重要的；否则，不同的进程可能会对应该是同一块的不同副本做出冲突的更改。每个缓存块存储在 {\tt struct buf} \fileref{kernel/buf.h} 中。{\tt struct buf} 有一个锁字段，有助于确保一次只有一个进程使用给定的磁盘块。然而，这个锁是不够的：如果块根本不在缓存中，两个进程想要同时使用它怎么办？没有 {\tt struct buf}（因为块尚未缓存），因此没有什么可以锁定的。Xv6 通过将额外的锁（{\tt bcache.lock}）与缓存块的身份集相关联来处理这种情况。需要检查块是否已缓存的代码（例如，{\tt bget} \lineref{kernel/bio.c:/^bget/}）或更改缓存块集的代码必须持有 {\tt bcache.lock}；在该代码找到块和它需要的 {\tt struct buf} 之后，它可以释放 {\tt bcache.lock} 并只锁定特定的块。这是一种常见模式：一个锁用于项集，加上每个项一个锁。

通常，获取锁的同一个函数会释放它。但更精确的看法是，锁在必须表现为原子的序列开始时获取，在该序列结束时释放。如果序列在不同的函数、不同的线程或不同的 CPU 上开始和结束，那么锁的获取和释放也必须这样做。锁的作用是强制其他使用者等待，而不是将某段数据固定到特定的执行者。一个例子是 {\tt yield} 中的 {\tt acquire}\lineref{kernel/proc.c:/^yield/}，它是在调度器线程中释放的，而不是在获取它的进程中。另一个例子是 {\tt ilock} 中的 {\tt acquiresleep}\lineref{kernel/fs.c:/^ilock/}；这段代码在读取磁盘时经常睡眠；它可能在不同的 CPU 上唤醒，这意味着锁可能在不同的 CPU 上获取和释放。

释放受嵌入在对象中的锁保护的对象是一项微妙的业务，因为拥有锁不足以保证释放是正确的。问题情况发生在其他线程在 {\tt acquire} 中等待使用对象时；释放对象隐式释放嵌入的锁，这将导致等待线程故障。一种解决方案是跟踪对对象的引用数量，以便仅在最后一个引用消失时才释放它。参见 {\tt pipeclose} \lineref{kernel/pipe.c:/^pipeclose/} 作为示例；{\tt pi->readopen} 和 {\tt pi->writeopen} 跟踪管道是否有文件描述符引用它。

通常，你会在对相关项目集进行读写序列周围看到锁；锁确保其他线程只看到已完成的更新序列（只要它们也加锁）。如果更新是对单个共享变量的简单写入呢？例如，\texttt{setkilled} 和 \texttt{killed} \lineref{kernel/proc.c:/^setkilled/} 在它们对 \lstinline{p->killed} 的简单使用周围加锁。如果没有锁，一个线程可能在另一个线程读取它的同时写入 \lstinline{p->killed}。这是一个 \indextextx{竞态（race condition）}，C 语言规范说竞态产生 \indextext{未定义行为（undefined behavior）}，这意味着程序可能崩溃或产生不正确的结果。锁防止竞态并避免未定义行为。

竞态可能破坏程序的一个原因是，如果没有锁或等效构造，编译器可能生成以与原始 C 代码截然不同的方式读写内存的机器代码。例如，调用 \texttt{killed} 的线程的机器代码可能将 \lstinline{p->killed} 复制到寄存器，并且只读取那个缓存值；这意味着该线程可能永远看不到对 \lstinline{p->killed} 的任何写入。锁防止这种缓存。

% sleep locks.
% hand-over-hand locking in namei.
%   namei's use of refcount to prevent changing underfoot,
%     and lock when actually using it.
% example of where deadlock is avoided? namei?
% spawn a thread to evade a lock order problem?

\section{类似锁的模式}

在许多地方，xv6 以类似锁的方式使用引用计数或标志来指示对象已分配且不应被释放或重新使用。进程的 {\tt p->state} 以这种方式工作，{\tt file}、{\tt inode} 和 {\tt buf} 结构中的引用计数也是如此。虽然在每种情况下锁都保护标志或引用计数，但正是后者防止对象被过早释放。

文件系统使用 {\tt struct inode} 引用计数作为一种可以被多个进程持有的共享锁，以避免如果代码使用普通锁会发生的死锁。例如，{\tt namex} \lineref{kernel/fs.c:/^namex/} 中的循环依次锁定路径名命名的每个目录。然而，{\tt namex} 必须在循环结束时释放每个锁，因为如果它持有多个锁，如果路径名包含点（例如，{\tt a/./b}），它可能会与自己死锁。它也可能与涉及目录和 {\tt ..} 的并发查找死锁。如 Chapter~\ref{CH:FS} 所解释的，解决方案是让循环将其引用计数递增的目录 inode 带到下一次迭代，但不锁定。

某些数据项在不同时间受不同机制保护，有时可能由 xv6 代码的结构隐式保护免受并发访问，而不是由显式锁保护。例如，当物理页面空闲时，它受 \texttt{kmem.lock} \lineref{kernel/kalloc.c:/^. kmem;/} 保护。如果页面然后被分配为管道 \lineref{kernel/pipe.c:/^pipealloc/}，它受不同的锁保护（嵌入的 \lstinline{pi->lock}）。如果页面被重新分配给新进程的用户内存，它根本不受锁保护。相反，分配器不会将该页面给予任何其他进程（直到它被释放）这一事实保护它免受并发访问。
新进程内存的所有权很复杂：首先在 {\tt fork} 中父进程分配和操作它，然后子进程使用它，（子进程退出后）父进程再次拥有内存并将其传递给 {\tt kfree}。这里有两个教训：数据对象在其生命周期的不同时间点可能以不同方式免受并发保护，保护可能采取隐式结构的形式而不是显式锁。

最后一个类似锁的例子是在调用 {\tt mycpu()} \lineref{kernel/proc.c:/^myproc/} 周围需要禁用中断。禁用中断使调用代码相对于可能强制上下文切换的定时器中断是原子的，从而将进程移动到不同的 CPU。

\section{完全没有锁}

有几个地方 xv6 完全没有锁就共享可变数据。一个是在自旋锁的实现中，尽管可以将 RISC-V 原子指令视为依赖于硬件中实现的锁。另一个是 {\tt main.c} 中的 {\tt started} 变量 \lineref{kernel/main.c:/^volatile/}，用于防止其他 CPU 运行，直到 CPU 零完成初始化 xv6；{\tt volatile} 确保编译器实际生成加载和存储指令。

Xv6 包含一个 CPU 或线程写入一些数据，另一个 CPU 或线程读取数据，但没有特定锁专门保护该数据的情况。例如，在 {\tt fork} 中，父进程写入子进程的用户内存页，子进程（不同的线程，可能在不同的 CPU 上）读取这些页面；没有锁显式保护这些页面。这严格来说不是锁问题，因为子进程在父进程完成写入后才开始执行。这是一个潜在的内存排序问题（参见 Chapter~\ref{CH:LOCK}），因为如果没有内存屏障，没有理由期望一个 CPU 看到另一个 CPU 的写入。然而，由于父进程释放锁，子进程在启动时获取锁，{\tt acquire} 和 {\tt release} 中的内存屏障确保子进程的 CPU 看到父进程的写入。

\section{并行性}

锁主要用于在正确性的利益下抑制并行性。因为性能也很重要，内核设计者通常必须考虑如何以既能实现正确性又能允许并行性的方式使用锁。虽然 xv6 不是为高性能系统设计的，但考虑哪些 xv6 操作可以并行执行，哪些可能在锁上冲突仍然值得。

Xv6 中的管道是适度良好并行性的一个例子。每个管道有自己的锁，因此不同进程可以在不同 CPU 上并行读写不同管道。对于给定管道，然而，写入者和读取者必须等待对方释放锁；它们不能同时读/写同一管道。从空管道读取（或写入满管道）必须阻塞的情况也是如此，但这不是由于锁定方案。

上下文切换是一个更复杂的例子。两个内核线程，每个在自己的 CPU 上执行，可以同时调用 {\tt yield}、{\tt sched} 和 {\tt swtch}，调用将并行执行。线程各自持有一个锁，但它们是不同的锁，因此它们不必等待对方。然而，一旦进入 {\tt scheduler}，两个 CPU 可能在搜索 {\tt RUNNABLE} 进程表时冲突在锁上。也就是说，xv6 可能在上下文切换期间从多个 CPU 获得性能提升，但也许不如它可能获得的那么多。

另一个例子是来自不同 CPU 上不同进程的 {\tt fork} 的并发调用。调用可能不得不等待彼此的 {\tt pid\_lock} 和 {\tt kmem\_lock}，以及在进程表中搜索 {\tt UNUSED} 进程所需的每个进程锁。另一方面，两个 fork 进程可以完全并行地复制用户内存页和格式化页表页。

上述每个例子中的锁定方案在某些情况下牺牲了并行性能。在每种情况下，可以使用更复杂的设计获得更多并行性。是否值得取决于细节：相关操作被调用的频率，代码在持有争用锁时花费多长时间，多少 CPU 可能同时运行冲突操作，代码的其他部分是否是更具限制性的瓶颈。很难猜测给定的锁定方案是否可能导致性能问题，或者新设计是否显著更好，因此通常需要在实际工作负载上进行测量。

\section{练习}

\begin{enumerate}

\item 修改 xv6 的管道实现，允许在不同 CPU 上同时读取和写入同一管道。

\item 修改 xv6 的 \texttt{scheduler()} 以减少当不同 CPU 同时查找可运行进程时的锁争用。

\item 消除 xv6 的 \texttt{fork()} 中的一些串行化。

\end{enumerate}
